\documentclass[10pt,a4paper]{amsart}
\usepackage[margin=19mm]{geometry}
\usepackage{amsmath,amssymb,mathtools}
\usepackage[scr=rsfs,cal=euler]{mathalfa}
\usepackage{xfrac}
\usepackage[unicode,hidelinks,bookmarks=false]{hyperref}
\newcommand{\ie}{i.e.\ }
\newcommand{\cf}{cf.\ }
\newcommand{\resp}{respectively\ }
\newcommand{\Subsection}{\subsection}
\providecommand{\nobreakdash}{\nobreak\hbox{-}}
\setlength{\emergencystretch}{2em}
\multlinegap0pt
\usepackage{mathtools}
\usepackage{stmaryrd}
\usepackage{amssymb}
\usepackage{enumerate}
\usepackage{tikz}
\usepackage[shortlabels]{enumitem}
\usepackage{caption}
\usepackage{booktabs}
\newtheorem{defn}{Definition}
\newtheorem{rmk}{Remark}
\newtheorem{lem}{Lemma}
\newtheorem{prop}{Proposition}
\title{\large \bf $\mathsf{GL}_N(\mathbb{C})$ BROWNIAN MOTION AND STOCHASTIC PDE ON ENTIRE FUNCTIONS}
\author{\small THEODOROS ASSIOTIS AND ZAHRA SADAT MIRSAJJADI}
\date{Source: arXiv:2605.06429v1 (CC BY 4.0)}
\begin{document}
\maketitle

\noindent\emph{Source and licence.} \large \bf $\mathsf{GL}_N(\mathbb{C})$ BROWNIAN MOTION AND STOCHASTIC PDE ON ENTIRE FUNCTIONS. Version arXiv:2605.06429v1 (CC BY 4.0), distributed by arXiv under the Creative Commons Attribution 4.0 International licence, http://creativecommons.org/licenses/by/4.0/, as recorded in the arXiv licence metadata for this exact version and shown on the abstract page https://arxiv.org/abs/2605.06429v1. Full licence notice: \texttt{licenses/arxiv-2605.06429-CC-BY-4.0.txt}.\par\smallskip

\noindent\textit{Editorial scope.} Counted unit: the article's prop PropositionLimitingGibbs (source0.tex lines 736--746). The article's complete original proof of it is delivered with it (source0.tex lines 748--859, one \texttt{proof} environment of 112 lines). No mathematics is written, repaired or re-derived: the statement and the proof are the article's own lines, in the article's own notation, and no retained line is substituted or re-flowed.  Retained uncounted as the source-local context the proof needs: lines 610--624; lines 616--618; lines 626--640; lines 692--709.  Nothing was omitted from any retained span, so the package carries no omission record.  No retained line contains a citation, so the wrapper carries no bibliography and no bibliography entry.  No retained line contains a figure environment, an image-inclusion command or drawing code, so no image is delivered and the package declares no asset dependency.  Editorial formatting only, in addition: an AMS \texttt{amsart} preamble that restates the article's own declarations and macros used by the retained lines, namely the environments \texttt{defn}, \texttt{rmk}, \texttt{lem}, \texttt{prop} on separate counters, together with the standard packages those lines need.  The retained environments print in this document as Definition 1, Remark 1, Lemma 1, Proposition 1 in order of appearance, whereas the article prints the same items with the numbering of its own full document.  The complete native source of the article as distributed in the arXiv e-print for this exact version is bundled unchanged as \texttt{sources/200020/source0.tex}.
\begin{defn}\label{GibbsDefinition}
Let $N \in \mathbb{N}\cup \{\infty\}$. We say that a line ensemble $\mathcal{L}\in \mathrm{C}(\llbracket 1, N \rrbracket)$ satisfies the Gibbs resampling property with respect to the law of the $\mathsf{A}$-bridge if for all $\llbracket \ell,\ell+k-1 \rrbracket \subset \llbracket 1, N \rrbracket$ with $\ell, k \ge 1$ and any $[a,b]\in \mathbb{R}_+$ we have
\begin{equation}\label{GibbsPropertyFormal}
\mathsf{Law}\left(\mathcal{L}|_{\llbracket \ell,\ell+k-1 \rrbracket \times [a,b]}|\mathscr{F}_{\mathrm{ext}}\left(\llbracket \ell, \ell+k-1 \rrbracket \times (a,b) \right)\right)=\mathbf{P}_{\mathrm{avoid}}^{a,b;\mathbf{x},\mathbf{y};u,l},\ \ \textnormal{ almost surely,}
\end{equation}
where the external sigma-algebra $\mathscr{F}_{\mathrm{ext}}\left(\llbracket \ell, \ell+k-1 \rrbracket \times (a,b) \right)$ is defined as
\begin{equation}\label{ExternalSigmaAlgebra}
\mathscr{F}_{\mathrm{ext}}\left(\llbracket \ell, \ell+k-1 \rrbracket \times (a,b) \right)= \sigma\left(\mathcal{L}_j(t);(j,t)\notin \llbracket \ell, \ell+k-1 \rrbracket\times (a,b)\right), 
\end{equation}
and the (random) boundary data $\mathbf{x},\mathbf{y},u,l$ given by,
\begin{align*}
\mathbf{x}=\left(\mathcal{L}_\ell(a),\mathcal{L}_{\ell+1}(a),\dots,\mathcal{L}_{k+\ell-1}(a)\right), \  \mathbf{y}=\left(\mathcal{L}_{\ell}(b),\mathcal{L}_{\ell+1}(b),\dots,\mathcal{L}_{k+\ell-1}(b)\right), \ u=\mathcal{L}_{\ell-1},  \ l=\mathcal{L}_{\ell+k},
\end{align*}
with the convention $\mathcal{L}_{0}\equiv \infty$, $\mathcal{L}_{N+1}\equiv -\infty$.
\end{defn}

\begin{equation}
\mathscr{F}_{\mathrm{ext}}\left(\llbracket \ell, \ell+k-1 \rrbracket \times (a,b) \right)= \sigma\left(\mathcal{L}_j(t);(j,t)\notin \llbracket \ell, \ell+k-1 \rrbracket\times (a,b)\right), 
\end{equation}

\begin{rmk}\label{RmkNonIntersectionGibbs}
Observe that, in our definition of the Gibbs property it is implicit that for all fixed $t\ge 0$, almost surely $(\mathcal{L}_i(t))_{i\in \llbracket 1,N \rrbracket }\in \mathbf{W}_N^\circ$ in order for the right hand side of \eqref{GibbsPropertyFormal} to be well-defined. In fact, a stronger statement holds: almost surely for all $t\ge 0$ we have $(\mathcal{L}_i(t))_{i\in \llbracket 1,N \rrbracket }\in \mathbf{W}_N^\circ$ as we spell out now. Write $\mathbf{P}$ and $\mathbf{E}$ for $\mathsf{Law}(\mathcal{L})$ and expectation with respect to it. It suffices to show that for all $k \in \mathbb{N}$ and times $T\ge 0$ we have,
\begin{equation*}
\mathbf{P}\left(\inf_{t\in [0,T]}\left|\mathcal{L}_k(t)-\mathcal{L}_{k+1}(t)\right|>0\right)=1.
\end{equation*}
From \eqref{GibbsPropertyFormal} we have for any bounded measurable function $\mathsf{F}:\mathrm{C}(\{k,k+1\}\times[0,T])\to \mathbb{R}_+$,
\begin{equation*}
\mathbf{E}\left[\mathsf{F}\left(\mathcal{L}|_{\{k,k+1\}\times[0,T]}\right)\big|\mathscr{F}_{\mathrm{ext}}\left(\{k,k+1\} \times (0,T) \right)\right]=\mathbf{E}_{\mathrm{avoid}}^{0,T;\mathbf{x},\mathbf{y};u,l}\left[\mathsf{F}(\mathsf{Z})\right], \ \textnormal{ almost surely},
\end{equation*}
where $\mathsf{Z}$ is distributed according to $\mathbf{P}_{\mathrm{avoid}}^{0,T;\mathbf{x},\mathbf{y};u,l}$, and $\mathbf{x},\mathbf{y},u,l$ as in Definition \ref{GibbsDefinition}. Taking expectations we get,
\begin{equation*}
\mathbf{E}\left[\mathsf{F}\left(\mathcal{L}|_{\{k,k+1\}\times[0,T]}\right)\right]=\mathbf{E}\left[\mathbf{E}_{\mathrm{avoid}}^{0,T;\mathbf{x},\mathbf{y};u,l}\left[\mathsf{F}(\mathsf{Z})\right]\right].
\end{equation*}
 Choosing $\mathsf{F}(f,g)=\mathbf{1}_{f(t)>g(t), \forall t \in [0,T]}$ we get the desired statement since the inner expectation is 1.
\end{rmk}

\begin{lem}\label{Lem-CondDistConv}
Assume that for all $[a,b] \subset \mathbb{R}_+$, and $x_N,y_N \in \mathbb{R}$ such that $x_N \to x$, $y_N \to y$ as $N \to \infty$, we have the weak convergence of probability measures,
\begin{equation}\label{SingleBridgeConv}
\mathbf{P}_N^{a,b;x_N,y_N} \longrightarrow \mathbf{P}^{a,b;x,y}, \ \textnormal{ as } N \longrightarrow \infty,
\end{equation}
where $\mathbf{P}_N^{a,b;x_N,y_N}$ corresponds to the law of a $\mathsf{A}_N$-bridge and $\mathbf{P}^{a,b;x,y}$ to the law of a $\mathsf{A}$-bridge.  Let $k\in \mathbb{N}$ and $\mathbf{x}_N,\mathbf{y}_N, \mathbf{x},\mathbf{y}\in \mathbf{W}_{k}^\circ$ with $\mathbf{x}_N \to \mathbf{x}, \mathbf{y}_N \to \mathbf{y}$ as $N \to \infty$. Moreover, suppose that $u_N,u \in \mathrm{C}([a,b];\mathbb{R}\cup \{ \infty \})$ and $l_N,l \in \mathrm{C}([a,b];\mathbb{R}\cup \{ -\infty \})$ are such that 
\begin{equation*}
u(t)>l(t), \ \forall t \in [a,b], \ u(a)>x_1, \ u(b)>y_1, \ l(a)<x_k, \ l(b)<y_k,
\end{equation*}
and also, as $N \to \infty$, uniformly in $[a,b]$,
\begin{equation}\label{BoundaryConvergence}
u_N \longrightarrow u, \ \ l_N \longrightarrow l.
\end{equation}
Then, for $N \in \mathbb{N}$ large enough, the conditional law $\mathbf{P}_{\mathrm{avoid};N}^{a,b;\mathbf{x}_N,\mathbf{y}_N;u_N,l_N}$ is well-defined and we have the following weak convergence of probability measures, as $N \longrightarrow \infty$,
\begin{equation}\label{ConditionedBridgeConv}
\mathbf{P}_{\mathrm{avoid};N}^{a,b;\mathbf{x}_N,\mathbf{y}_N;u_N,l_N}\longrightarrow \mathbf{P}_{\mathrm{avoid}}^{a,b;\mathbf{x},\mathbf{y};u,l}.
\end{equation}
\end{lem}

\begin{prop}\label{PropositionLimitingGibbs}
Let $(\mathcal{L}^{(N)})_{N\in\mathbb{N}} $ be a sequence of line ensembles such that for each $N\in\mathbb{N}$, with $\mathcal{L}^{(N)}\in \mathrm{C}(\llbracket 1, N \rrbracket \times \mathbb{R}_+)$  satisfies the Gibbs resampling property with respect to the law of a $\mathsf{A}_N$-bridge. Let $\mathcal{L}$ be a $\mathrm{C}(\mathbb{N}\times\mathbb{R}_+)$-valued line ensemble. Moreover, assume that as $N \longrightarrow \infty$, for all $M \in \mathbb{N}$ (when $N\ge M$),
\begin{align*} \left(\mathcal{L}^{(N)}_{i}(t)\right)_{i\in \llbracket 1, M \rrbracket, t\in \mathbb{R}_+}\overset{\mathrm{d}}{\longrightarrow}\left(\mathcal{L}_{i}(t)\right)_{i\in \llbracket 1, M \rrbracket, t\in \mathbb{R}_+} \ \textnormal{ in }\; \mathrm{C}\left(\llbracket 1, M \rrbracket \times\mathbb{R}_+\right), 
\end{align*}
and that for any fixed $t\ge 0$, almost surely $(\mathcal{L}_i(t))_{i\in \mathbb{N}}\in \mathbf{W}_\infty^\circ$. Finally, we assume that for any $[a,b]\subset \mathbb{R}_+$ and any $x_N,y_N \in \mathbb{R}$, whenever $x_N \to x$ and $y_N \to y$, we have the weak convergence of probability measures
\begin{equation*}
\mathbf{P}_N^{a,b;x_N,y_N} \longrightarrow\mathbf{P}^{a,b;x,y}, \ \ \textnormal{ as } N \longrightarrow\infty,
\end{equation*}
where $\mathbf{P}_N^{a,b;x_N,y_N}$ and $\mathbf{P}^{a,b;x,y}$ correspond to the law of a $\mathsf{A}_N$-bridge and $\mathsf{A}$-bridge respectively. Then, $\mathcal{L}$ satisfies the Gibbs resampling property with respect to the law of an $\mathsf{A}$-bridge.
In particular, the paths of $\mathcal{L}$ are almost surely non-intersecting.
\end{prop}

\begin{proof}

We fix an index set
$ \llbracket \ell , \ell+k-1 \rrbracket \subset \mathbb{N}$, with $\ell,k\ge 1$,
and an arbitrary time interval \([a,b]\subset{\mathbb{R}_+}\).
Recall the definition of the external sigma-algebra $\mathscr{F}_{\mathrm{ext}}\left(\llbracket \ell, \ell+k-1 \rrbracket \times (a,b) \right)$ from \eqref{ExternalSigmaAlgebra}. To prove the result, we need to show that
for any bounded Borel-measurable function
$\mathsf{F}: \mathrm{C}\left(\llbracket \ell, \ell+k-1 \rrbracket \times [a,b]\right) \to \mathbb{R}_+$,
we have, with $\mathbf{E}$ denoting expectation with respect to $\mathsf{Law}(\mathcal{L})$,
\begin{equation}\label{GibbsGoalIdentity}
\mathbf{E}\left[ \mathsf{F}\left(\mathcal{L}|_{\llbracket \ell, \ell+k-1 \rrbracket \times [a,b]}\right) \big| \mathscr{F}_{\mathrm{ext}}\left(\llbracket \ell, \ell+k-1\rrbracket\times (a,b)\right) \right]
= \mathbf{E}_{\mathrm{avoid}}^{\mathbf{x},\mathbf{y};u,l}\left[\mathsf{F}(\mathsf{Z})\right], \  \ \textnormal{almost surely},  
\end{equation}
where $\mathsf{Z}$ has the law
$\mathbf{P}_{\textnormal{avoid}}^{\mathbf{x},\mathbf{y};u,l}$,
and, the random boundary data is given by
\begin{align*}
\mathbf{x}&=\left(\mathcal{L}_{\ell}(a),\mathcal{L}_{\ell+1}(a),\dots,\mathcal{L}_{\ell+k-1}(a)\right),\  \ \mathbf{y}=
\left(\mathcal{L}_{\ell}(b),\mathcal{L}_{\ell+1}(b),\dots,\mathcal{L}_{\ell+k-1}(b)\right),
\\ u&=
\mathcal{L}_{\ell-1}\  \  (u=+\infty, \textnormal{ if } \ell=1),\hspace{7mm}\  \ l=
\mathcal{L}_{\ell+k}.
\end{align*}
Observe that the right hand side of \eqref{GibbsGoalIdentity} is well-defined almost surely by virtue of the assumption that for any fixed $t\ge 0$, almost surely $(\mathcal{L}_i(t))_{i\in \mathbb{N}}\in \mathbf{W}_\infty^\circ$. Here and below, we omit the parameters $a,b$ from the notation for brevity.


By the definition and since
$ \mathbf{E}_{\mathrm{avoid}}^{\mathbf{x}, \mathbf{y}; u, l}\left[\mathsf{F}(\mathsf{Z})\right]
$
is  $\mathscr{F}_{\mathrm{ext}}\left(\llbracket \ell, \ell+k-1\rrbracket \times (a,b)\right)$-measurable,
it suffices to  verify that
\begin{align*}
\mathbf{E}\left[\mathbf{1}_{\mathscr{B}}\mathsf{F}\left(\mathcal{L}|_{\llbracket \ell, \ell+k-1 \rrbracket \times [a,b]}\right)\right] = \mathbf{E}\left[\mathbf{1}_{\mathscr{B}}\mathbf{E}_{\mathrm{avoid}}^{\mathbf{x},\mathbf{y};u,l}\left[\mathsf{F}(\mathsf{Z})\right]\right],
    \  \ \forall \mathscr{B}\in\mathscr{F}_{\mathrm{ext}}\left(\llbracket \ell, \ell+k-1\rrbracket \times (a,b)\right).
\end{align*}
Consider now the following $\pi$-system of cylinder sets,
\begin{align*}
\mathfrak{C} \overset{\textnormal{def}}{=}  \bigg\{ \mathscr{B}&=
\left\{\left (\mathcal{L}_{i_1}(t_1), \ldots, \mathcal{L}_{i_m}(t_m)\right) \in A \right\}: \\
&(i_j,t_j) \notin \llbracket \ell, \ell+k-1 \rrbracket \times (a,b),\,j\in\llbracket 1,m\rrbracket,\,
m\in\mathbb{N},\; A \subseteq \mathbb{R}^m \textnormal{ Borel}\bigg\}.
\end{align*}
From standard measure-theoretic results, we have 
\begin{align*}
\sigma\left(\mathfrak{C}\right)=\mathscr{F}_{\mathrm{ext}}\left(\llbracket \ell, \ell+k-1 \rrbracket \times (a,b)\right).
\end{align*}
Next, define
\begin{align*}
\mathfrak{D} \overset{\textnormal{def}}{=}  \left\{ \mathscr{B}\in \mathscr{F}_{\mathrm{ext}}\left(\llbracket \ell, \ell+k-1 \rrbracket \times (a,b)\right): \mathbf{E}\left[\mathbf{1}_{\mathscr{B}}\mathsf{F}\left(\mathcal{L}|_{\llbracket \ell, \ell+k-1\rrbracket \times [a,b]}\right)\right] = \mathbf{E}\left[\mathbf{1}_{\mathscr{B}}\, \mathbf{E}_{\mathrm{avoid}}^{\mathbf{x},\mathbf{y};u,l}\left[\mathsf{F}(\mathsf{Z})\right]\right] \right\}.
\end{align*}
We show in the following that
\begin{align}\label{ConditionalExpectationIdentity}
\mathbf{E}\left[\mathbf{1}_{\mathscr{B}}\mathsf{F}\left(\mathcal{L}|_{\llbracket \ell, \ell+k-1 \rrbracket \times [a,b]}\right)\right] = \mathbf{E}\left[\mathbf{1}_{\mathscr{B}}\mathbf{E}_{\mathrm{avoid}}^{\mathbf{x},\mathbf{y};u,l}\left[\mathsf{F}(\mathsf{Z})\right]\right],
    \  \ \forall \mathscr{B}\in\mathfrak{C}.
\end{align}
Assuming \eqref{ConditionalExpectationIdentity}, it is straightforward to verify, using the monotone convergence theorem, that \(\mathfrak{D}\) is a (Dynkin) \(\lambda\)-system containing $\mathfrak{C}$. Hence, the desired result follows by
applying the %(Dynkin) 
$\pi-\lambda$ theorem.

We now proceed to proving the identity
\eqref{ConditionalExpectationIdentity}.  To establish \eqref{ConditionalExpectationIdentity}, we begin by considering the case where $\mathsf{F}$ is continuous (and bounded). By a density argument it moreover suffices to  prove the following. For any $m \in \mathbb{N}$,  $(i_j,t_j) \notin \llbracket \ell, \ell+k-1 \rrbracket \times (a,b),\,j\in\llbracket 1,m\rrbracket$ and any continuous bounded function $\mathfrak{H}$ from $\mathbb{R}^m$ to $\mathbb{R}$ we have:
\begin{align}\label{ConditionalExpectWithFunction}
&\mathbf{E}\left[\mathfrak{H}\left (\mathcal{L}_{i_1}(t_1), \ldots, \mathcal{L}_{i_m}(t_m)\right)\mathsf{F}\left(\mathcal{L}|_{\llbracket \ell, \ell+k-1 \rrbracket \times [a,b]}\right)\right]\nonumber \\&= \mathbf{E}\left[\mathfrak{H}\left (\mathcal{L}_{i_1}(t_1), \ldots, \mathcal{L}_{i_m}(t_m)\right)\mathbf{E}_{\mathrm{avoid}}^{\mathbf{x},\mathbf{y};u,l}\left[\mathsf{F}(\mathsf{Z})\right]\right].
\end{align}
By the Skorokhod representation theorem, we realize $(\mathcal{L}^{(N)})_{N\in \mathbb{N}}$ and $\mathcal{L}$ on a common probability space $(\Omega,\mathbf{Q})$ such that, $\mathbf{Q}$-almost surely, as $N \longrightarrow \infty$, for all $M \in \mathbb{N}$ (with $N \ge M$),
\begin{align}\label{a.s.conv} \left(\mathcal{L}^{(N)}_{i}(t)\right)_{i\in \llbracket 1, M \rrbracket, t\in \mathbb{R}_+}\longrightarrow\left(\mathcal{L}_{i}(t)\right)_{i\in \llbracket 1, M \rrbracket, t\in \mathbb{R}_+} \ \textnormal{ in }\; \mathrm{C}\left(\llbracket 1, M \rrbracket \times\mathbb{R}_+\right).
\end{align}
Hence, $\mathbf{Q}$-almost surely, as $N \longrightarrow \infty$ we have,
\begin{equation}\label{ExtFnConv}
\mathfrak{H}\left (\mathcal{L}^{(N)}_{i_1}(t_1), \ldots, \mathcal{L}^{(N)}_{i_m}(t_m)\right) \longrightarrow \mathfrak{H}\left (\mathcal{L}_{i_1}(t_1), \ldots, \mathcal{L}_{i_m}(t_m)\right).
\end{equation}
Moreover, for each \(N\in\mathbb{N}\)
with $N\ge  \ell+k$, by virtue of the Gibbs property of $\mathcal{L}^{(N)}$ one obtains with $\mathfrak{H}$ as above
\begin{align}
&\mathbf{Q}\left[\mathfrak{H}\left (\mathcal{L}^{(N)}_{i_1}(t_1), \ldots, \mathcal{L}^{(N)}_{i_m}(t_m)\right)\mathsf{F}\left(\mathcal{L}^{(N)}|_{\llbracket \ell, \ell+k-1 \rrbracket \times [a,b]}\right)\right]\nonumber
\\&= \mathbf{Q}\left[\mathfrak{H}\left (\mathcal{L}^{(N)}_{i_1}(t_1), \ldots, \mathcal{L}^{(N)}_{i_m}(t_m)\right) \mathbf{E}_{N,\mathrm{avoid}}^{\mathbf{x}^{(N)},\mathbf{y}^{(N)};u^{(N)},l^{(N)}}\left[\mathsf{F}\left(\mathsf{Z}^{(N)}\right)\right]\right]\label{BGP_N},
\end{align}
where
$\mathsf{Z}^{(N)}$
is 
$\mathbf{P}_{N,\mathrm{avoid}}^{\mathbf{x}^{(N)},\mathbf{y}^{(N)};u^{(N)},l^{(N)}}$-distributed and the random boundary data is given by,
\begin{align*}
\mathbf{x}^{(N)}&=
\left(\mathcal{L}_{\ell}^{(N)}(a),\mathcal{L}_{\ell +1}^{(N)}(a),\dots,\mathcal{L}_{\ell+k-1}^{(N)}(a)\right),\  \ \mathbf{y}^{(N)}=
\left(\mathcal{L}_{\ell}^{(N)}(b),\mathcal{L}_{\ell+1}^{(N)}(b),\dots,\mathcal{L}_{\ell+k-1}^{(N)}(b)\right),
\\ u^{(N)}&=
\mathcal{L}_{\ell-1}^{(N)},\  \  (u^{(N)}=+\infty, \textnormal{ if } \ell=1),\hspace{2mm}\  \ l^{(N)}=
\mathcal{L}_{\ell+k}^{(N)}.
\end{align*} 
In light of \eqref{a.s.conv}
and \eqref{ExtFnConv}, 
we have using the bounded convergence theorem  and moreover using Lemma \ref{Lem-CondDistConv} (recall that for any fixed $t\ge 0$, almost surely $(\mathcal{L}_i(t))_{i\in \mathbb{N}}\in \mathbf{W}_\infty^\circ$)
that
\begin{align}
&\mathbf{Q}\left[\mathfrak{H}\left (\mathcal{L}^{(N)}_{i_1}(t_1), \ldots, \mathcal{L}^{(N)}_{i_m}(t_m)\right) \mathsf{F}\left(\mathcal{L}^{(N)}|_{\llbracket \ell, \ell+k-1\rrbracket \times [a,b]}\right)\right]\nonumber \\
&\xrightarrow{N \to \infty}
\mathbf{Q}\left[\mathfrak{H}\left (\mathcal{L}_{i_1}(t_1), \ldots, \mathcal{L}_{i_m}(t_m)\right)\mathsf{F}\left(\mathcal{L}|_{\llbracket \ell, \ell+k-1\rrbracket \times [a,b]}\right)\right], \label{RHS-conv}
\\
&\mathbf{Q}\left[\mathfrak{H}\left (\mathcal{L}^{(N)}_{i_1}(t_1), \ldots, \mathcal{L}^{(N)}_{i_m}(t_m)\right)\mathbf{E}_{N,\mathrm{avoid}}^{\mathbf{x}^{(N)},\mathbf{y}^{(N)},u^{(N)},l^{(N)}}\left[\mathsf{F}\left(\mathsf{Z}^{(N)}\right)\right]\right]\nonumber
\\
&\xrightarrow{N \to \infty}
\mathbf{Q}\left[\mathfrak{H}\left (\mathcal{L}_{i_1}(t_1), \ldots, \mathcal{L}_{i_m}(t_m)\right)\mathbf{E}_{\mathrm{avoid}}^{\mathbf{x},\mathbf{y};u,l}[\mathsf{F}(\mathsf{Z})]\right].\label{LHS-conv}
\end{align}
Therefore, \eqref{ConditionalExpectWithFunction} and thus \eqref{ConditionalExpectationIdentity} with $\mathsf{F}$ continuous and bounded follows from
\eqref{RHS-conv} and \eqref{LHS-conv} by taking $N\to\infty$ in \eqref{BGP_N}.
Now, let us define
\begin{align*}
    \mathscr{G}\overset{\textnormal{def}}{=}  \bigg\{ &\mathsf{F} :\mathrm{C}\left(\llbracket \ell, \ell+k-1\rrbracket \times [a,b]\right)\to\mathbb{R}_+:  \\
    &\mathsf{F} \textnormal{ is a bounded, Borel-measurable function satisfying }\eqref{ConditionalExpectationIdentity}  \bigg\}.
\end{align*}
It is straightforward to verify that \(\mathscr{G}\) is a monotone class. Moreover, from the argument above, \(\mathscr{G}\) contains the algebra of bounded continuous functions on \(\mathrm{C}\left(\llbracket \ell, \ell+k-1\rrbracket \times [a,b]\right)\), which generates the Borel \(\sigma\)-algebra on this space. Therefore, by the monotone class theorem, \(\mathscr{G}\) contains all bounded Borel-measurable functions on \(\mathrm{C}\left(\llbracket \ell, \ell+k-1\rrbracket \times [a,b]\right)\) which completes the proof of \eqref{ConditionalExpectationIdentity}. The very final statement then follows from Remark \ref{RmkNonIntersectionGibbs}.
\end{proof}

\end{document}
